% GENERATED FILE — DO NOT EDIT; authoritative prose is Markdown.
% source: companion/dynamic-orthing-noetic-learning-and-orthability.md
% source-sha256: e626ad00c5eb947a2295f653538266f3b87c8b7732a7070a50273780fdd31f0b
% source-qualification: research-stage-draft
% source-qualification: not-peer-reviewed
% source-qualification: not-externally-peer-reviewed
% source-qualification: not-empirically-validated
% source-qualification: candidate-terminology-not-adopted
% source-qualification: philosophical-conclusions-conditional-on-stated-premises
% source-qualification: engineering-evidence-does-not-support-metaphysical-claims
\documentclass[10pt,letterpaper,twocolumn]{article}
\usepackage{amsmath}
\usepackage{amssymb}
\usepackage{booktabs}
\usepackage{geometry}
\usepackage{hyperref}
\usepackage{microtype}
\usepackage{natbib}
\usepackage{xcolor}
\geometry{letterpaper,margin=0.75in}
\hypersetup{colorlinks=true,linkcolor=blue,urlcolor=blue,citecolor=blue}
\setlength{\emergencystretch}{3em}
\title{Dynamic orthing, noetic learning, and dynamic orthability}
\author{}
\date{}
\begin{document}
\twocolumn[
\maketitle
\textbf{Status:} DRAFT companion — philosophical placement, conditional on stated premises. OPUS CANDIDATE — REQUIRES FRESH FABLE REVIEW BEFORE MERGE. Creed- internal material is school-labeled (Atharī) where it appears. \textbf{No empirical claim; no theological claim is asserted as established; asserts no soul access.}

\vspace{1em}
]

\section*{1. Source-status note}
Five source statuses are used and never conflated: \textbf{primary-text verified} (a located classical text), \textbf{secondary reconstruction} (a modern scholarly compilation or dissertation), \textbf{cross-source synthesis} (this project's joining of sources), \textbf{orthemological extension} (a claim of this framework), and \textbf{creed-internal inference} (a step valid only inside a labeled school, e.g. the Atharī route on rungs 8–10). Every load-bearing claim below and every rung in \href{\detokenize{DYNAMIC-ORTHABILITY-ARGUMENT-MAP.yaml}}{\texttt{DYNAMIC-\allowbreak{}ORTHABILITY-\allowbreak{}ARGUMENT-\allowbreak{}MAP.\allowbreak{}yaml}} carries one of these five labels; a creed-internal step is never presented as school-neutral, and no empirical result is ever labeled above cross-source synthesis. The supplied \texttt{Ibn Taymiyyah's Epistemology.md} is a \textbf{collection} of modern scholarly materials, not one authorial work by Ibn Taymiyya; the "concrete reason / ideal reason" contrast is a \textbf{modern interpretive lens} (associated with Evans and applied by Turner), not Ibn Taymiyya's own fixed technical pair. El-Tobgui's study of the \emph{Darʾ Taʿāruḍ al-ʿaql wa-l-naql} is a \textbf{secondary reconstruction}. Nothing here converts a secondary compilation into primary-text verification; doctrinal loci carry a verification queue in the sourcing ledger.


\section*{2. Bearer ontology (the distinctions that must not collapse)}

\begin{itemize}
\item \textbf{faculty / disposition} — the capacity to judge (creed-internally, fiṭrah as a normative readiness, not a measured faculty);

\item \textbf{represented standard} $\tilde\mu_{\alpha,t}$ — an adopted/embodied governing rule;

\item \textbf{metaortheme} $\mu$ — the repeatable normative type;

\item \textbf{metaorthemma} $\bar\mu_{e,j}$ — the case binding;

\item \textbf{execution} — what was done;

\item \textbf{placement / judgment} $\hat p$ — the inferred profile (belief, never ground truth);

\item \textbf{objective result} $O^*(m; A)$ — the actual profile under the analysis;

\item \textbf{truth linkage} — whether the placement is truth-linked.

\end{itemize}
These are the epistemological bearers a dynamic account must keep separate (Decisions 0025, 0027, 0028). Orthemology's contribution is the \textbf{typed separation}; it does not adjudicate the scholarship's own claims.


\section*{3. Concrete versus sound reason}
"Concrete" describes \emph{instantiation} or historical embodiment; "sound" describes a declared \emph{normative/epistemic status}. Neither entails the other: concrete reasoning can be sound or unsound. Soundness here is \textbf{claim-relative} — its proper bearer and dimensions are those of StrictlySoundReasoning, i.e. $\operatorname{StrictlySoundReasoning}_q(e) :=\allowbreak{} \operatorname{ReasoningPathAdequate}_q(e) \wedge\allowbreak{} \operatorname{TOKEN\_TRUTH\_LINKED}_q(e)$ (Decision 0011), not a free-floating predicate. The Taymiyyan reconstruction (secondary) locates the load-bearing contrast within concretely instantiated reason — \emph{ʿaql ṣarīḥ} (clear reason through a sound disposition and valid sources) versus reason corrupted by inherited premises, whim, conjecture, habit, motive, or imitation — not between "the concrete" and "the sound."


\section*{4. Fiṭrah and proper function}
Within the creed-internal (Atharī) model, fiṭrah is a \textbf{normative disposition / proper-function orientation}, not one measured faculty, one coordinate, one metaortheme, one algorithm, or one quantity inferred from discourse (Decision 0025 §4; fixture N19 keeps runtime closure separate from restoration). Any "minimum-entropy attractor" phrasing is provisional and creed-internal.


\section*{5. Tawātur versus memetic propagation}
Meta-noetic memetics makes the false-tawātur problem central (Decision 0028; fixtures TW1–TW6). Many apparent reports may be independent corroboration, copies of one upstream source, repetitions produced by one institutional rule, mutations of one carrier, or convergence after genuinely independent routes. The Taymiyyan material treats tawātur as depending on \textbf{more than number} — origin, path independence, dependence, and qualitative/circumstantial indicants. The ecology $\Gamma^\mu$ supplies exactly this structure; \textbf{count and graph degree establish no warrant}.


\section*{6. Mental / external anti-reification}
Latent states, noetic fields, orthemes, metaorthemes, and orthability are \textbf{formal / modelling / analysis objects}. Objectivity lies in the truth of the analysis-relative instantiation/placement relation constrained by concrete reality ($\operatorname{Inst}_A$, $O^*(m; A)$ once $A$ is fixed — Decision 0001; symbol normalization is Decision 0005), \textbf{not} in separately subsisting Platonic universals (Decision 0009 §6).


\section*{7. Multi-target noetic orthing}
A discourse occurrence $m_{\text{disc}}$ is \textbf{evidence about} a possible subject target $m_{\text{subj}}$, never identical to it (Decision 0027). A deformation state-type, if placed, belongs to the subject under a declared analysis — never to the utterance because the utterance is observed. Every noetic claim names its target, evidence, scope, analysis, status, and non-claims; no operation manipulates a soul-state.


\section*{8. Dynamic orthability (candidate concept)}
The metaphysical companion treats \textbf{static} conditions of evaluability. Dynamic orthing raises the \textbf{dynamic} question: what must reality be like for error to be discoverable, distinctions learnable, correction trajectories possible, and governing standards revisable — \textbf{without truth being created by the learner}?

This sits inside \textbf{Decision 0010's three-sense apparatus} (orthability-L/O/R) and adds a temporal/trajectory question; it introduces no new sense. \textbf{Candidate:} dynamic orthability is the analysis-relative availability of lawful, evidence- sensitive, correction-capable trajectories among unresolved, mistaken, and more adequate states of representation and action under a declared analysis $A$.


\subsection*{8.1 Eight modalities that must not collapse}
"Correction-capable trajectories are available" is ambiguous across \textbf{eight} distinct modalities (audit B39); collapsing any two is a fallacy:


\begin{enumerate}
\item \textbf{metaphysical possibility} — a correction trajectory exists in some possible world;

\item \textbf{nomological possibility} — one exists under the actual laws;

\item \textbf{objective reachability} — $m' \in \operatorname{Reach}_A(m, \Pi)$: some admissible action-sequence reaches it (a fact about the state space, actor-independent);

\item \textbf{controllability} — a \emph{governed} policy can steer the system to it (reachability plus a control law), still actor-independent;

\item \textbf{actor accessibility} — $m' \in \operatorname{Accessible}_{A,\alpha,t}(m, H_t)$: \emph{this} actor can reach it from \emph{its} history and capacities;

\item \textbf{evidence-relative discoverability} — it is findable under the evidence actually available to the inquiry;

\item \textbf{representation-class learnability} — this representation/hypothesis class \emph{can} learn the correction structure at all;

\item \textbf{convergence} — a learning procedure is \emph{guaranteed} to converge to it (rate/limit claims), a strictly stronger claim than learnability.

\end{enumerate}
Dynamic orthability is an \textbf{objective corrigibility} notion at levels 1–4 (the correction structure and its reachability/controllability are not created by any learner); it must \textbf{not} be collapsed into 5–8 (agent-relative accessibility, evidence-relative discoverability, representation-class learnability, guaranteed convergence). Levels 3–4 are \emph{objective but modal}, not "already actual"; only levels 1–2 are unconditioned. Dynamic orthability does \textbf{not} entail guaranteed learnability, universal convergence, one scalar gradient (Decision 0025/0033 hold G1 as proposed-candidate, not literal G2), a created learner producing the modal order (orthability-O is not created, Decision 0010), or the metaphysical bridge.


\subsection*{8.2 Two numbered scales — do not conflate}
This paper uses \textbf{two} numbered scales. The \textbf{eight modalities} (§8.1) disambiguate the \emph{meaning} of "a correction is available"; the \textbf{fourteen argument nodes} (§9, \texttt{DYNAMIC-\allowbreak{}ORTHABILITY-\allowbreak{}ARGUMENT-\allowbreak{}MAP.\allowbreak{}yaml}) are the metaphysical ladder. They are different scales. Modalities 1–4 (objective corrigibility) are what rung 2 (dynamic corrigibility) asserts; modalities 5–8 (agent-relative accessibility, discoverability, learnability, convergence) are what rung 3 (agent-accessible learnability) concerns and what OSM exemplifies. "Level \emph{n}" in §8.1 therefore never means "rung \emph{n}". Dynamic orthability is defined at modalities 1–4 and addressed by rungs 1–3; neither scale's lower entries license any upper entry of the other.


\section*{9. Metaphysical argument map}
Each rung is a \textbf{separate premise}, with an inference type, an objection, and a rival exit. OSM and DAEE are \textbf{worked explananda / applications} at the lower levels, not proof of the upper bridge.


\subsection*{Generated typed argument-node summary}
Argument nodes: \textbf{14}.


% breakable-row-table: data-rows=14 total-rendered-characters=2195 max-row-rendered-characters=211
\par\medskip
\hrule height 0.8pt
\smallskip
% breakable-row: 1/14
\noindent\textbf{Order}: 1\par
\noindent\textbf{Stable ID}: \texttt{ARG-01}\par
\noindent\textbf{Node}: local evaluability\par
\noindent\textbf{Scope}: cross-framework-dialectical\par
\noindent\textbf{Inference / bridge}: analytic / bounded\par
\noindent\textbf{Claim role}: orthemological-extension\par
\noindent\textbf{Dependency}: none\par
\noindent\textbf{Rival exit}: deflationism\par
\smallskip
\hrule height 0.4pt
\smallskip
% breakable-row: 2/14
\noindent\textbf{Order}: 2\par
\noindent\textbf{Stable ID}: \texttt{ARG-02}\par
\noindent\textbf{Node}: objective corrigibility\par
\noindent\textbf{Scope}: cross-framework-dialectical\par
\noindent\textbf{Inference / bridge}: explanatory / conditional\par
\noindent\textbf{Claim role}: orthemological-extension\par
\noindent\textbf{Dependency}: ARG-01\par
\noindent\textbf{Rival exit}: humean-regularity\par
\smallskip
\hrule height 0.4pt
\smallskip
% breakable-row: 3/14
\noindent\textbf{Order}: 3\par
\noindent\textbf{Stable ID}: \texttt{ARG-03}\par
\noindent\textbf{Node}: bounded representational adaptation\par
\noindent\textbf{Scope}: cross-framework-dialectical\par
\noindent\textbf{Inference / bridge}: computational-analogy / illustrative-only\par
\noindent\textbf{Claim role}: computational-analogy\par
\noindent\textbf{Dependency}: ARG-02\par
\noindent\textbf{Rival exit}: instrumentalism\par
\smallskip
\hrule height 0.4pt
\smallskip
% breakable-row: 4/14
\noindent\textbf{Order}: 4\par
\noindent\textbf{Stable ID}: \texttt{ARG-04}\par
\noindent\textbf{Node}: normative proper-function bridge\par
\noindent\textbf{Scope}: cross-framework-dialectical\par
\noindent\textbf{Inference / bridge}: teleological / conditional\par
\noindent\textbf{Claim role}: cross-source-synthesis\par
\noindent\textbf{Dependency}: ARG-02, ARG-03\par
\noindent\textbf{Rival exit}: selected-function-naturalism\par
\smallskip
\hrule height 0.4pt
\smallskip
% breakable-row: 5/14
\noindent\textbf{Order}: 5\par
\noindent\textbf{Stable ID}: \texttt{ARG-05}\par
\noindent\textbf{Node}: underived locus\par
\noindent\textbf{Scope}: cross-framework-dialectical\par
\noindent\textbf{Inference / bridge}: transcendental-scope-limited / conditional\par
\noindent\textbf{Claim role}: orthemological-extension\par
\noindent\textbf{Dependency}: ARG-01\par
\noindent\textbf{Rival exit}: aseity-bootstrapping\par
\smallskip
\hrule height 0.4pt
\smallskip
% breakable-row: 6/14
\noindent\textbf{Order}: 6\par
\noindent\textbf{Stable ID}: \texttt{ARG-06}\par
\noindent\textbf{Node}: intellectual ground\par
\noindent\textbf{Scope}: cross-framework-dialectical\par
\noindent\textbf{Inference / bridge}: metaphysical-explanatory / conditional\par
\noindent\textbf{Claim role}: cross-source-synthesis\par
\noindent\textbf{Dependency}: ARG-05\par
\noindent\textbf{Rival exit}: impersonal-ground\par
\smallskip
\hrule height 0.4pt
\smallskip
% breakable-row: 7/14
\noindent\textbf{Order}: 7\par
\noindent\textbf{Stable ID}: \texttt{ARG-07}\par
\noindent\textbf{Node}: modal status of the ground\par
\noindent\textbf{Scope}: cross-framework-dialectical\par
\noindent\textbf{Inference / bridge}: modal / held\par
\noindent\textbf{Claim role}: cross-source-synthesis\par
\noindent\textbf{Dependency}: ARG-05, ARG-06\par
\noindent\textbf{Rival exit}: modal-primitivism\par
\smallskip
\hrule height 0.4pt
\smallskip
% breakable-row: 8/14
\noindent\textbf{Order}: 8\par
\noindent\textbf{Stable ID}: \texttt{ARG-08}\par
\noindent\textbf{Node}: will\par
\noindent\textbf{Scope}: cross-framework-dialectical\par
\noindent\textbf{Inference / bridge}: metaphysical / conditional\par
\noindent\textbf{Claim role}: cross-source-synthesis\par
\noindent\textbf{Dependency}: ARG-06\par
\noindent\textbf{Rival exit}: brute-contingency\par
\smallskip
\hrule height 0.4pt
\smallskip
% breakable-row: 9/14
\noindent\textbf{Order}: 9\par
\noindent\textbf{Stable ID}: \texttt{ARG-09}\par
\noindent\textbf{Node}: power\par
\noindent\textbf{Scope}: cross-framework-dialectical\par
\noindent\textbf{Inference / bridge}: metaphysical / conditional\par
\noindent\textbf{Claim role}: cross-source-synthesis\par
\noindent\textbf{Dependency}: ARG-08\par
\noindent\textbf{Rival exit}: impersonal-efficacy\par
\smallskip
\hrule height 0.4pt
\smallskip
% breakable-row: 10/14
\noindent\textbf{Order}: 10\par
\noindent\textbf{Stable ID}: \texttt{ARG-10}\par
\noindent\textbf{Node}: Wisdom and fittingness\par
\noindent\textbf{Scope}: athari-taymiyyan-operative\par
\noindent\textbf{Inference / bridge}: creed-internal-teleological / held-pending-source-custody\par
\noindent\textbf{Claim role}: creed-internal-inference\par
\noindent\textbf{Dependency}: ARG-04, ARG-06\par
\noindent\textbf{Rival exit}: euthyphro-fittingness\par
\smallskip
\hrule height 0.4pt
\smallskip
% breakable-row: 11/14
\noindent\textbf{Order}: 11\par
\noindent\textbf{Stable ID}: \texttt{ARG-11}\par
\noindent\textbf{Node}: capacity for disclosure\par
\noindent\textbf{Scope}: cross-framework-dialectical\par
\noindent\textbf{Inference / bridge}: metaphysical / conditional\par
\noindent\textbf{Claim role}: cross-source-synthesis\par
\noindent\textbf{Dependency}: ARG-06\par
\noindent\textbf{Rival exit}: capacity-scepticism\par
\smallskip
\hrule height 0.4pt
\smallskip
% breakable-row: 12/14
\noindent\textbf{Order}: 12\par
\noindent\textbf{Stable ID}: \texttt{ARG-12}\par
\noindent\textbf{Node}: actual divine Speech\par
\noindent\textbf{Scope}: athari-taymiyyan-operative\par
\noindent\textbf{Inference / bridge}: revelational-school-internal / revelational\par
\noindent\textbf{Claim role}: creed-internal-inference\par
\noindent\textbf{Dependency}: ARG-11\par
\noindent\textbf{Rival exit}: alternative-kalam-articulations\par
\smallskip
\hrule height 0.4pt
\smallskip
% breakable-row: 13/14
\noindent\textbf{Order}: 13\par
\noindent\textbf{Stable ID}: \texttt{ARG-13}\par
\noindent\textbf{Node}: Speech bearer classification\par
\noindent\textbf{Scope}: athari-taymiyyan-operative\par
\noindent\textbf{Inference / bridge}: typed-creed-internal-distinction / bounded\par
\noindent\textbf{Claim role}: creed-internal-inference\par
\noindent\textbf{Dependency}: ARG-12\par
\noindent\textbf{Rival exit}: alternative-kalam-articulations\par
\smallskip
\hrule height 0.4pt
\smallskip
% breakable-row: 14/14
\noindent\textbf{Order}: 14\par
\noindent\textbf{Stable ID}: \texttt{ARG-14}\par
\noindent\textbf{Node}: Rabb lexical disambiguation candidate\par
\noindent\textbf{Scope}: source-bounded-proposed-crosswalk\par
\noindent\textbf{Inference / bridge}: orthemological-extension / held-pending-individually-verified-lexical-sources\par
\noindent\textbf{Claim role}: orthemological-extension\par
\noindent\textbf{Dependency}: none\par
\noindent\textbf{Rival exit}: source-custody-hold\par
\smallskip
\hrule height 0.8pt
\par\medskip
The table is generated from the structured map. Cross-framework scope means dialectical accessibility and rival routing, not a neutral or coequal tribunal. The Atharī/Taymiyyan frame is the declared operative frame. The common-premise fittingness-to-Wisdom bridge remains held; actual divine Speech remains revelational and school-internal; OSM and DAEE validate neither metaphysics nor theology.

\textbf{Gap invariance (rung 3 → rung 4).} The step from agent-accessible learnability (rung 3) to normative proper function (rung 4) is the framework's critical inductive gap — a regularity → norm (broadly is/ought) transition. It is \textbf{invariant under additional exemplification}: no quantity of further OSM/DAEE-type learning results, however many cases converge, can close it, because more instances of what a faculty \emph{does} do not by themselves yield a fact about what it \emph{ought} to do. Rung 4 is carried only by the explanatory/teleological bridge premise (rival exit: teleosemantic naturalism), never by learning evidence. Empirical density at rung 3 therefore leaves rung 4's status exactly where it was.

Rungs 1–3 are what the static and dynamic Orthemology material addresses. Rungs 4–10 each require \textbf{explicit bridge premises already identified in the companion papers}; the learning/diagnosis material never carries them. Rung 10 is creed-internal (Atharī) and school-labeled. The full machine-readable map — every rung's premise, inference type, conclusion, dependency, evidence/source status, strongest objection, rival exit, and school-neutral vs Atharī-internal status — is \href{\detokenize{DYNAMIC-ORTHABILITY-ARGUMENT-MAP.yaml}}{\texttt{DYNAMIC-\allowbreak{}ORTHABILITY-\allowbreak{}ARGUMENT-\allowbreak{}MAP.\allowbreak{}yaml}}, checked in CI. Its \texttt{non\_entailments} block records that OSM/DAEE exemplify only rungs 1–3 and prove no upper rung.


\section*{10. Theological / Atharī route and the created/uncreated distinction}
Rung 10 (revelation-specific Speech) is developed in full in the Atharī companion (\texttt{orthability-\allowbreak{}divine-\allowbreak{}attributes-\allowbreak{}and-\allowbreak{}speech-\allowbreak{}athari.\allowbreak{}md}), under explicit school-labeled, revelational premises. The reason/revelation firewall stands: \textbf{no engineering or learning evidence supports any theological rung}, and OSM/DAEE belong only to the lower explanatory levels.

This section adds only a \textbf{bounded dynamic cross-section}, not a replacement doctrine. The dynamic/learning vocabulary of this paper must keep the following strictly separate (none is reducible to another, and none is established by OSM or DAEE):


\begin{itemize}
\item \textbf{created human linguistic convention} — the historically evolved conventions of a human language (a created, dynamic, learned system);

\item \textbf{created token and media} — the created sound, ink, page, screen, and recitation through which wording is borne;

\item \textbf{eternally known linguistic possibility} — that a given wording is possible/known is not itself a created event;

\item \textbf{the divine act of speaking} — a divine act, not a creaturely learning trajectory;

\item \textbf{revealed Arabic wording} — the specific revealed wording, received in a created human language;

\item \textbf{creaturely reception} — the created hearing, recitation, and writing by which creatures receive it;

\item \textbf{the Qur'an as divine Speech} — under the Atharī route, divine Speech (uncreated), distinguished from the created media that convey it.

\end{itemize}
The dynamic-orthing apparatus of this paper models only the \textbf{created, learnable} side (convention, media, reception) and the analysis-relative structure of intelligibility (orthability-L/O/R). It says nothing for or against the uncreated status of divine Speech; that is a creed-internal claim carried solely by the Atharī companion's labeled premises. Collapsing "created Arabic media are learned/dynamic" into "divine Speech is created" is exactly the conflation the firewall forbids.


\section*{11. Objections and rival exits}

\begin{itemize}
\item \textbf{"OSM proves normative proper function."} No — OSM empirically exemplifies learning dynamics (rung 3); it does not establish rung 4. Reply: keep the explanandum/proof distinction (row 3 vs 4).

\item \textbf{"DAEE's restorative norms prove their own ground."} No — DAEE \emph{presupposes} the norms it applies (epistemological circularity); it is a worked explanandum, not a proof (Decision 0025; fixtures N5/N17).

\item \textbf{"Dynamic orthability entails a converging learner."} No — it is levels 1–3, explicitly not 4–6.

\end{itemize}

\section*{12. Limitations}
This is a placement, not a proof. The bridge premises (rungs 4–10) are argued elsewhere and remain conditional. Primary-text verification of the Taymiyyan loci is a standing queue. No empirical study validates any claim here.

\textbf{Absence conditions (when dynamic orthability fails under A).} The candidate is falsifiable-in-principle relative to a declared analysis $A$. Dynamic orthability is \emph{absent} under $A$ when any of: (i) placements lack determinate correctness conditions, so rung 1 fails and there is nothing to be corrected toward; (ii) no more-adequate state is reachable — $\operatorname{Reach}_A(m, \Pi) = \varnothing$ for every unresolved or mistaken $m$ — so the correction structure is empty; or (iii) apparent improvement is shown to be brute regularity with no truth-linkage (rung 2's rival exit obtains and "correction" tracks no objective gradient). Stating these keeps the concept from being unfalsifiable; showing a real case meets (i)–(iii) is itself analysis-relative and empirical, not settled here.


\section*{13. References}

\begin{itemize}
\item Sun, W., et al. (2025). \emph{Learning produces an orthogonalized state machine in the hippocampus.} \textbf{Nature} 640. DOI 10.1038/s41586-024-08548-w (CC BY 4.0).

\item El-Tobgui, C. S. (2013). \emph{Reason, Revelation and the Reconstitution of Rationality: Ibn Taymiyya's Darʾ Taʿāruḍ al-ʿAql wa-l-Naql.} PhD dissertation, McGill University (secondary reconstruction).

\item Turner, J. B. (2022). \emph{Ibn Taymiyya on Theistic Signs and Knowledge of God.} \textbf{Religious Studies} 58 (secondary).

\item Evans, C. S. (2010). \emph{Natural Signs and Knowledge of God.} Oxford University Press (secondary; the "natural signs" lens).

\item Companion papers: \texttt{orthability-\allowbreak{}and-\allowbreak{}the-\allowbreak{}ground-\allowbreak{}of-\allowbreak{}intelligibility.\allowbreak{}md}, \texttt{orthability-\allowbreak{}divine-\allowbreak{}attributes-\allowbreak{}and-\allowbreak{}speech-\allowbreak{}athari.\allowbreak{}md}; Decisions 0010, 0024, 0025, 0027, 0028.

\end{itemize}

\section*{Cross-references (bounded)}
Dynamic orthing / latent learning: \texttt{applications/\allowbreak{}latent-\allowbreak{}state-\allowbreak{}orthing/\allowbreak{}} (Decision 0024). Meta-noetic memetics / corrective dynamics: \texttt{applications/\allowbreak{}daee-\allowbreak{}epistemics/\allowbreak{}} (Decision 0025). Multi-target noetic orthing: Decision 0027. Represented standards + ecology: Decision 0028.


\twocolumn
\bibliographystyle{plainnat}
\bibliography{../../../references/orthemology}
\end{document}
